\documentclass[a4paper,12pt]{article}
\usepackage{amsmath}
\usepackage{graphicx}
\usepackage{multicol}
\usepackage{geometry}
\usepackage{fancyhdr}
\usepackage{tikz}
\usepackage{eso-pic}

\geometry{top=1.5cm, bottom=1.5cm, left=1.5cm, right=1.5cm}

\pagestyle{fancy}
\fancyhf{}
\rhead{Esc. 4-117 Ejército de los Andes}
\lhead{Termodinámica y Máquinas Térmicas}
\cfoot{\thepage}

\AddToShipoutPicture*{
    \put(150,550){\rotatebox{45}{\textcolor[gray]{0.95}{EXAMEN ORIGINAL NO DIGITALIZAR}}}
}

\begin{document}

\begin{center}
    \Large \textbf{Evaluación de Termodinámica y Máquinas Térmicas} \\
    \normalsize Duración: 60 minutos \\
    Alumno/a: \underline{\hspace{6cm}} \hspace{0.5cm} Fecha: \underline{\hspace{2cm}}
\end{center}

\vspace{0.3cm}

\begin{multicols}{2}

\section*{Problemas}

\begin{enumerate}
\item (\textbf{Ley de Charles})\\
Un gas ocupa $1.8\,L$ a $280\,K$. ¿Cuál será su volumen a $350\,K$ si la presión permanece constante?

\item (\textbf{Ley de Gay-Lussac})\\
Un gas está a $1.0\,atm$ y $10^\circ C$. ¿Qué presión tendrá a $90^\circ C$ si el volumen no cambia?

\item (\textbf{Ecuación general de los gases ideales})\\
Un gas ocupa $3.0\,L$ a $2.0\,atm$ y $300\,K$. ¿Cuál será su volumen si la presión baja a $1.2\,atm$ y la temperatura sube a $330\,K$?

\item (\textbf{Gas ideal – cálculo de volumen})\\
¿Cuál es el volumen que ocupan $0.5\,mol$ de gas a $1\,atm$ y $273\,K$? ($R = 0.082\,\frac{L \cdot atm}{mol \cdot K}$)

\item (\textbf{Ley de Dalton – presiones parciales})\\
Una mezcla gaseosa contiene $0.4\,mol$ de oxígeno, $0.3\,mol$ de nitrógeno y $0.3\,mol$ de hidrógeno. Si la presión total es $1.2\,atm$, calcular la presión parcial de cada gas.

\item (\textbf{Ley de Amagat – volúmenes parciales})\\
Tres gases ocupan individualmente los siguientes volúmenes a la misma temperatura y presión: $V_O = 2.0\,L$, $V_N = 1.5\,L$, $V_H = 2.5\,L$. ¿Cuál es el volumen total de la mezcla? ¿Qué fracción del volumen total ocupa cada gas?

\end{enumerate}

\vspace{0.2cm}

\end{multicols}

\end{document}
